\subsection{Что такое Python и зачем он нужен?}

Python — это высокоуровневый язык программирования, который был создан Гвидо ван Россумом и впервые выпущен в 1991 году. Назван он не в честь змеи, а в честь популярного британского комедийного шоу <<Летающий цирк Монти Пайтона>>.

\subsubsection*{Почему Python?}
\begin{itemize}
    \item \textbf{Простота:} Синтаксис Python похож на обычный английский язык.
    \item \textbf{Универсальность:} Используется в веб-разработке, науке, AI и робототехнике.
    \item \textbf{Мощь:} Огромное количество библиотек (OpenCV, NumPy, TensorFlow) для решения сложных задач.
\end{itemize}

Для программирования квадрокоптеров Pioneer мы выбрали Python, так как он позволяет легко взаимодействовать с оборудованием и обрабатывать видеопоток с камеры.

\subsubsection*{Сравнение Python и Lua}
Мы уже знакомы с языком Lua, который используется для написания скриптов автопилота. Давайте посмотрим, чем Python отличается от Lua:

\begin{table}[H]
\centering
\begin{tabularx}{\textwidth}{>{\bfseries}l X X}
\toprule
Характеристика & \textbf{Lua} & \textbf{Python} \\
\midrule
Назначение & Встраиваемый скриптовый язык (для игр, контроллеров) & Универсальный язык общего назначения \\
Синтаксис & Использует \texttt{end} для блоков кода & Использует \textbf{отступы} для блоков кода \\
Индексация & С 1 (первый элемент — 1) & С 0 (первый элемент — 0) \\
Библиотеки & Ограниченный набор & Огромная экосистема (PyPI) \\
Скорость & Очень быстрый & Медленнее, но мощнее \\
\bottomrule
\end{tabularx}
\caption{Сравнение Lua и Python}
\end{table}

\begin{infobox}[Важное отличие]
Главное, что нужно запомнить: в Python блоки кода (циклы, функции, условия) выделяются \textbf{отступами} (обычно 4 пробела), а не ключевыми словами \texttt{do ... end}.
\end{infobox}

\subsubsection*{Первая программа}
Традиционно изучение любого языка начинается с программы, которая выводит приветствие на экран.

\begin{codefile}[python]{python/01_python_basics/01_hello_world.py}
\end{codefile}

Здесь функция \texttt{print()} выводит текст, заключенный в кавычки, в консоль.

\begin{taskbox}[Привет, мир!]
Напишите программу, которая выводит ваше имя и фамилию на экран.
\end{taskbox}
